\mychapter{II. LITERATURE REVIEW}{II. LITERATURE REVIEW}

\section{Introduction to Academic and Methodological Foundations}
\hs This chapter presents a systematic review of the academic literature, central bank empirical research, and international statistical standards relevant to high-frequency inflation measurement. Tracking consumer price dynamics using alternative digital data sources represents one of the most transformative innovations in applied economic measurement. The review is structured around five interconnected thematic pillars that directly support the design and execution of the Cambodia Daily Consumer Price Index (CPI) Medallion Pipeline:
\begin{enumerate}
    \item \textbf{The Transition from Physical Price Surveys to Automated Big Data Ingestion:} The theoretical and operational shift from traditional in-person enumerator surveys to automated daily web collection systems.
    \item \textbf{Natural Language Processing and Artificial Intelligence in Product Classification:} Modern semantic representations, vector similarity methods, and supervised classification frameworks for mapping raw retail listings into standard international consumption taxonomies.
    \item \textbf{Econometric Index Number Theory and Mathematical Axioms:} Elementary price relatives, axiomatic comparisons between geometric and arithmetic aggregation, missing price imputation principles, and the elimination of calculation bias.
    \item \textbf{High-Frequency Inflation Forecasting via Machine Learning:} Methodological foundations for forward-looking consumer price prediction using high-frequency daily price signals, gradient boosted tree architectures, and non-linear lag dynamics.
    \item \textbf{Emerging Market Context and Research Gap:} Synthesizing findings from global implementations and highlighting the critical research gap in dollarized developing economies such as Cambodia.
\end{enumerate}

\section{From Manual Store Audits to Automated Digital Economic Sensing}
\subsection{The Historical Framework of Consumer Price Collection}
\hs For over a century, the standard international methodology for measuring consumer inflation has relied on periodic field surveys conducted by national statistical agencies. Under this traditional paradigm, trained field enumerators visit a fixed sample of brick-and-mortar retail outlets, physical markets, and service establishments on a monthly or quarterly schedule. Enumerators manually record price quotes on paper forms or handheld electronic tablets for a predetermined basket of goods and services.

\hs While this approach has provided a standardized foundation for national accounting worldwide, economists and statistical authorities have long recognized its inherent operational bottlenecks. Physical data collection is labor-intensive, logistically expensive, and constrained by sample size limitations. Field surveyors can only inspect a limited number of items during each store visit. More critically, the multi-week process required to gather, transmit, verify, clean, and aggregate field reports introduces a substantial publication lag. In most countries, official Consumer Price Index figures are released between 20 and 45 days after the end of the reference month. Under stable macroeconomic conditions, this lag is manageable; however, during episodes of rapid inflation, supply chain disruptions, currency volatility, or international commodity shocks, this delay leaves monetary and fiscal authorities operating with outdated information.

\subsection{The MIT Billion Prices Project and the Global Evidence for Online Prices}
\hs The modern era of alternative digital price measurement began with the groundbreaking work of the MIT Billion Prices Project \citep{cavallo2016billion}. Recognizing that retail businesses increasingly display full product catalogs and transaction prices on public e-commerce websites, Cavallo and Rigobon developed automated software crawlers to collect hundreds of thousands of daily price quotes across 22 countries.

\hs Their empirical findings established three fundamental principles that form the foundation of high-frequency price measurement:
\begin{itemize}
    \item \textbf{High Synchronization with Official Statistics:} In low-inflation and open economies (including the United States, Germany, Japan, and Brazil), price indices computed from daily web-scraped data demonstrated a Pearson correlation exceeding $0.98$ with official monthly CPI series compiled by national statistical agencies. This confirmed that prices published on digital retail platforms accurately reflect broad consumer price movements in the broader economy.
    \item \textbf{Early Warning Capabilities:} Because online prices are collected continuously every morning, web-scraped indices were able to detect macroeconomic turning points, retail promotional cycles, and foreign exchange pass-through shocks between 3 and 14 days after their occurrence, compared to the 30 to 60-day delay characteristic of traditional monthly surveys.
    \item \textbf{Statistical Truth and Anomaly Detection:} In countries experiencing severe inflation or official statistical distortions (such as Argentina during the 2007 to 2015 period), the Billion Prices Project demonstrated that online price crawling accurately tracked actual annual inflation of approximately $20.0\%$, while official statistical publications reported only $8.0\%$. This proved the objective value of automated public web data as an independent, tamper-proof source of macroeconomic intelligence.
\end{itemize}

\subsection{Institutional Adoption and Practical Web Scraping Standards}
\hs Following the success of academic initiatives, major national statistical institutes and international bodies began actively evaluating the incorporation of web scraping into official statistical production. Recognizing the growing importance of digital data, the Statistical Office of the European Communities formalized operational standards for statistical authorities \citep{eurostat2022webscraping}. The Eurostat framework outlines best practices for web scraping in the production of consumer price indices, emphasizing legal compliance, automated quality validation, robust crawler scheduling, and the need for resilient data engineering architectures capable of managing volatile online product catalogs.

\section{Natural Language Processing and Artificial Intelligence in Product Classification}
\subsection{The Challenge of Multilingual and Unstructured Product Data}
\hs A central challenge in transforming web-scraped data into an official Consumer Price Index is the absence of standardized product classifications in online retail stores. In traditional field surveys, human enumerators know exactly which predetermined category each item belongs to before recording its price. In contrast, web scrapers collect raw text titles, promotional descriptions, and unstructured attributes directly from website pages.

\hs In developing economies like Cambodia, this challenge is intensified by several distinct market characteristics:
\begin{itemize}
    \item \textbf{Multilingual Catalogues:} Online listings in Cambodia frequently mix Khmer script, English terminology, and French loan words within the same product title (for example, combining local food terms like \textit{Trey Chhpin} with international brand names and metric size units).
    \item \textbf{Lack of Universal Barcodes:} Unlike point-of-sale scanner databases from large department stores, online marketplace listings often lack standardized Global Trade Item Numbers (GTIN) or barcode identifiers. Products are identified primarily through free-form descriptive text.
    \item \textbf{Frequent Title and Formatting Variations:} Retailers regularly alter product titles to highlight temporary promotions, modify brand spelling, or include marketing adjectives (such as ``Special Deal'' or ``Fresh Arrival''), which causes traditional keyword searches and exact string matching to fail.
\end{itemize}

\subsection{Evolution from Rule-Based Matching to Transformer Language Models}
\hs Early attempts at automated product classification relied on keyword dictionaries and regular expression rules. While simple to implement, rule-based systems are brittle, difficult to maintain, and incapable of generalizing to new product varieties or unconventional phrasing.

\hs Recent advances in natural language processing (NLP) have transformed this domain through the use of deep transformer models and dense semantic vector representations. \cite{berki2025nlp} demonstrated that fine-tuning Bidirectional Encoder Representations from Transformers (BERT) on scraped e-commerce titles achieved a classification accuracy of \textbf{94.56\%} and a weighted precision of \textbf{94.07\%} when mapping raw product titles into 5-digit United Nations COICOP categories. The transformer architecture captures contextual relationships between words, enabling the system to distinguish between homonyms and interpret domain-specific jargon without manual human intervention.

\subsection{Central Bank Innovations: The BIS Project Spectrum}
\hs At the international central banking level, the Bank for International Settlements (BIS) Innovation Hub, in partnership with the European Central Bank (ECB) and Deutsche Bundesbank, launched \textbf{Project Spectrum} \citep{bis2024spectrum}. Project Spectrum investigated the application of modern Large Language Models (LLMs) and high-dimensional semantic text embeddings to classify billions of web-scraped price observations for real-time central bank inflation monitoring.

\hs The findings of Project Spectrum highlighted several key architectural principles:
\begin{itemize}
    \item Dense vector embeddings allow computers to represent products as points in a multi-dimensional mathematical space, where items with similar economic purposes naturally cluster together regardless of exact wording.
    \item Combining vector similarity with caching mechanisms (memoization) allows the system to remember previous classifications, reducing expensive computational calls by over $90\%$ after the initial database warm-up.
    \item Zero-shot and few-shot artificial intelligence prompts achieved classification precision exceeding \textbf{95\%} across multilingual retail catalogs, proving that modern language models can replace months of manual rule-crafting with automated, highly consistent categorization.
\end{itemize}

\subsection{Human-in-the-Loop Operational Architectures}
\hs Despite the high accuracy of modern artificial intelligence and machine learning models, national statistical integrity requires robust fallback mechanisms to handle ambiguous product descriptions, rare marketplace items, or low-confidence predictions. In high-frequency operational pipelines, this is achieved through Human-in-the-Loop (HITL) architectural design. By establishing statistical confidence thresholds, the pipeline automatically processes high-confidence classifications while flagging borderline items for expert review or deterministic rule overrides. This ensures that the automated system operates with maximum throughput without compromising the econometric consistency required for official price measurement.

\section{Econometric Index Number Theory and Axiomatic Foundations}
\subsection{The Elementary Level: Why Mathematical Formulas Matter}
\hs The mathematical aggregation of individual price observations into a national price index occurs in two distinct tiers: the \textbf{elementary level} (aggregating individual item quotes within an unweighted stratum) and the \textbf{macro level} (aggregating strata using national expenditure weights).

\hs As codified in the international gold standard \textit{Consumer Price Index Manual: Concepts and Methods} \citep{imf2020cpi}, published jointly by the IMF, ILO, OECD, Eurostat, UN, and World Bank, the choice of mathematical formula at the elementary level has profound consequences for the accuracy of inflation measurement. Historically, statistical agencies evaluated three primary elementary formulas:
\begin{enumerate}
    \item \textbf{The Carli Formula (Arithmetic Average of Price Relatives):} Calculates the simple arithmetic mean of price ratios between the current period and the base period:
    \begin{equation}
        P_{\text{Carli}}^{0:t} = \frac{1}{n} \sum_{i=1}^n \left( \frac{p_{i,t}}{p_{i,0}} \right)
    \end{equation}
    \item \textbf{The Dutot Formula (Ratio of Arithmetic Averages):} Calculates the ratio of the average price in the current period to the average price in the base period:
    \begin{equation}
        P_{\text{Dutot}}^{0:t} = \frac{\frac{1}{n} \sum_{i=1}^n p_{i,t}}{\frac{1}{n} \sum_{i=1}^n p_{i,0}}
    \end{equation}
    \item \textbf{The Jevons Formula (Geometric Mean of Price Relatives):} Calculates the unweighted geometric mean of price ratios (or equivalently, the ratio of geometric average prices):
    \begin{equation}
        P_{\text{Jevons}}^{0:t} = \prod_{i=1}^n \left( \frac{p_{i,t}}{p_{i,0}} \right)^{\frac{1}{n}} = \frac{\left( \prod_{i=1}^n p_{i,t} \right)^{\frac{1}{n}}}{\left( \prod_{i=1}^n p_{i,0} \right)^{\frac{1}{n}}}
    \end{equation}
\end{enumerate}

\subsection{Axiomatic Properties and the Elimination of Formula Drift}
\hs Econometric index theory evaluates formulas using mathematical axioms to ensure that an index behaves logically under all market conditions:
\begin{itemize}
    \item \textbf{The Time-Reversal Test:} If prices move from period 0 to period $t$, and then reverse exactly back to their starting levels in period 0, the product of the two price indices must equal exactly 1.0 ($I^{0:t} \times I^{t:0} = 1$). The Jevons formula satisfies this test unconditionally. In contrast, the Carli formula systematically fails the time-reversal test due to the mathematical inequality between arithmetic and harmonic means (Jensen's Inequality). If an item's price doubles from \$1 to \$2 and then drops back to \$1, the Carli formula calculates an average index of 1.25, erroneously reporting a $25\%$ increase in the cost of living even though prices did not change. In high-frequency daily price data, where retail sales and temporary promotional discounts cause prices to bounce frequently, the Carli formula accumulates an artificial upward drift of \textbf{1.0\% to 1.5\% per year}.
    \item \textbf{The Transitivity / Circularity Test:} An index comparing period 0 to period 2 should equal the chained product of indices from period 0 to period 1 and period 1 to period 2 ($I^{0:2} = I^{0:1} \times I^{1:2}$). The Jevons geometric formula satisfies circularity, ensuring that daily price indices can be chained over time without accumulating cumulative distortion.
    \item \textbf{Dimensional Invariance / Scale Independence:} An index must produce the same result regardless of whether prices are quoted in dollars, riels, cents, or thousands of riels. Both the Jevons and Carli formulas satisfy this test, whereas the Dutot formula is sensitive to the price scale of different items within the same category.
\end{itemize}

\hs Because the Jevons formula satisfies all fundamental axiomatic tests and models consumer substitution behavior under unitary elasticity, the IMF CPI Manual \citep{imf2020cpi} explicitly recommends the Jevons geometric mean as the standard elementary aggregator for unweighted price quotes.

\subsection{Macro-Level Aggregation and National Expenditure Weights}
\hs Once elementary price relatives are calculated for each product variety using the axiomatic Jevons formula, they must be combined into division-level indices and the final national Headline CPI. In online web scraping environments, transaction volumes and physical sales quantities for individual online listings are unobservable. Consequently, superlative index numbers requiring simultaneous high-frequency price and quantity data cannot be computed at the elementary level.

\hs Following the standard international methodology established by the IMF CPI Manual \citep{imf2020cpi}, high-frequency online systems bridge this gap through a two-tier aggregation architecture: elementary strata are compiled using unweighted geometric Jevons indices, while macro-level aggregation applies fixed expenditure weights derived from periodic national household surveys. In Cambodia, these weights are officially established by the Cambodia Socio-Economic Survey (CSES 2023), conducted by the National Institute of Statistics (NIS). This ensures that the high-frequency daily price series accurately reflects household expenditure patterns across all 12 UN COICOP divisions while remaining mathematically immune to formula drift.

\section{High-Frequency Inflation Forecasting via Machine Learning}

\subsection{The Micro-Price Revolution: Real-Time Baskets and Frequency Advantages}
\hs In macroeconomic forecasting, official inflation series have historically been modeled using low-frequency monthly or quarterly econometric specifications, such as Vector Autoregressions (VAR), Autoregressive Integrated Moving Average (ARIMA), and structural Phillips Curves. However, low-frequency models suffer from severe temporal inertia. In dynamic developing economies, by the time monthly official survey reports are compiled and published (typically 30 to 60 days after the reference period), supply chain shocks, fuel price shifts, or exchange rate fluctuations may have already propagated through the domestic retail economy.

\hs The advent of automated daily web scraping has catalyzed a fundamental paradigm shift toward \textbf{high-frequency micro-price forecasting} \citep{cavallo2016billion}. Rather than waiting for retrospective monthly aggregations, researchers can harvest continuous daily price trajectories across thousands of retail items. In their foundational work with the MIT Billion Prices Project, \cite{cavallo2016billion} demonstrated that daily web-scraped price series across 22 countries exhibited a Pearson correlation exceeding $0.98$ with official monthly CPI series in open economies, while detecting macroeconomic turning points and international commodity shocks 3 to 14 days after their occurrence.

\hs Extending this methodology to targeted consumption baskets, \cite{aparicio2021targeted} demonstrated that high-frequency web-scraped price microdata captures retail price stickiness, promotional markdowns, and intra-month expenditure shifts weeks before traditional statistical agencies publish monthly figures. Their empirical analysis showed that tracking targeted supermarket baskets in real time provides central banks with continuous early-warning surveillance, transforming inflation estimation from an episodic historical audit into an active, high-frequency predictive system.

\subsection{Non-Linear Machine Learning and the Superiority of Tree Ensembles}
\hs When modeling inflation from high-frequency multi-division price feeds, traditional linear regressions face severe theoretical and empirical limitations. Economic relationships in retail consumer markets are inherently non-linear and subject to structural regime shifts.

\hs In a comprehensive benchmark evaluation, \cite{medeiros2021forecasting} compared traditional econometric specifications against modern machine learning architectures (including LASSO, Ridge regression, Elastic Net, Factor Models, Random Forests, and Gradient Boosted Regression Trees) across multiple forecasting horizons. Their empirical findings revealed that \textbf{Gradient Boosted Regression Trees (GBRT) and Random Forests systematically outperformed traditional econometric benchmarks}, achieving a \textbf{10\% to 25\% reduction in Mean Squared Forecast Error (MSFE)} over standard autoregressions and Phillips Curve baselines. They established that non-linear algorithms excel because macroeconomic price formation is governed by complex threshold effects that linear shrinkage models cannot capture.

\hs In an extensive theoretical inquiry into why machine learning methods excel in macroeconomic prediction, \cite{gouletcoulombe2022machine} identified three structural advantages that explain why tree-based ensembles systematically dominate linear specifications:
\begin{itemize}
    \item \textbf{Asymmetric Input-Cost Pass-Through (``Rockets and Feathers'' Pricing):} Consumer prices frequently respond asymmetrically to input cost shocks. For example, retail fuel price surges trigger rapid, sharp upward price increases in transport and distribution (Division 07), whereas price declines transmit at a significantly slower rate. Decision trees naturally isolate these non-linear thresholds through recursive binary partitioning.
    \item \textbf{Automatic Multi-Way Interactions:} Tree ensembles discover complex interactions between concurrent macroeconomic shocks without requiring explicit polynomial feature engineering. In retail markets, a cost shock occurring during a major cultural celebration (such as Khmer New Year or Pchum Ben) exerts a compounding effect on final prices that linear additive terms fail to capture.
    \item \textbf{Robustness against Collinear High-Frequency Regressors:} In daily time-series feature stores containing dense multi-day autoregressive lags ($L_1, L_2, \dots, L_{30}$), linear models suffer from severe multicollinearity and variance inflation. Gradient boosted trees select the single most informative feature split at each node, maintaining robust generalizability across correlated high-frequency regressors.
\end{itemize}

\subsection{High-Frequency Food Price Dynamics and Central Bank Leading Indicators}
\hs In emerging and developing economies, aggregate inflation is predominantly led by supply-side shocks in agriculture and retail food distribution. In an extensive empirical study, \cite{macias2023nowcasting} analyzed over 2.4 million daily supermarket price quotes harvested across major retail chains in Poland. By integrating daily price microdata into recursive time-series and bridge forecasting models, their high-frequency food tracking framework achieved a \textbf{15\% to 35\% reduction in Root Mean Squared Forecast Error (RMSFE)} compared to traditional statistical baselines.

\hs The findings of \cite{macias2023nowcasting} provide two critical insights for high-frequency inflation forecasting:
\begin{enumerate}
    \item \textbf{Food as the Primary Leading Driver:} Food prices adjust with significantly higher daily frequency than sticky categories (such as housing or education), making high-frequency food price momentum the single strongest leading indicator for overall headline inflation.
    \item \textbf{Intra-Month Predictive Power:} The majority of the predictive gain from high-frequency scraping is realized within the first 10 to 14 days of the reference month. Tracking retail price arrivals during the first half of the month enables statistical analysts to project final monthly inflation with high precision weeks before official statistical releases.
\end{enumerate}

\subsection{Time-Series Econometrics, Validation Rigor, and Look-Ahead Discipline}
\hs Integrating high-frequency daily price series into predictive machine learning models requires strict methodological discipline to avoid data leakage and spurious correlation. In standard cross-sectional machine learning, randomized $k$-fold cross-validation is ubiquitous. However, in macroeconomic time series, random shuffling violates temporal ordering and leaks future price information into past training folds.

\hs \cite{babii2022machine} addressed this challenge by formalizing high-frequency time-series machine learning regressions under strict expanding-window (\textbf{walk-forward}) cross-validation protocols. By training exclusively on historical data points $1, \dots, t$ and evaluating out-of-sample performance strictly on forward horizons $t+1, \dots, t+H$, machine learning models provide genuine, unbiased out-of-sample generalization metrics (Root Mean Squared Error and Mean Absolute Error). Furthermore, direct multi-horizon target formulation avoids iterative compounding errors, ensuring stable projections across weekly, bi-weekly, and monthly horizons.

\subsection{Architectural Synthesis and Methodological Blueprint for Cambodia}
\hs The theoretical and empirical findings from these landmark studies establish the econometric and computational foundation for the Cambodia Daily CPI Medallion Pipeline. From \cite{cavallo2016billion} and \cite{aparicio2021targeted}, the pipeline adopts the principle of high-frequency daily retail price harvesting to bypass the standard 30 to 60-day official survey delays. From \cite{medeiros2021forecasting} and \cite{gouletcoulombe2022machine}, the pipeline deploys non-linear tree-based ensembles (LightGBM) to capture asymmetric pass-through and holiday interaction effects across dense autoregressive lags. From \cite{macias2023nowcasting}, the pipeline leverages high-frequency food price momentum (COICOP Division 01) as the primary leading indicator for aggregate consumer price fluctuations. Finally, following \cite{babii2022machine}, the forecasting engine enforces strict walk-forward cross-validation (\texttt{TimeSeriesSplit}) and direct multi-horizon targets ($H \in \{7, 14, 30\}$) to eliminate look-ahead bias and avoid compounding recursive forecast errors.

\section{The Emerging Market Context and Research Gap}
\hs While the integration of web scraping, artificial intelligence, and machine learning inflation forecasting has advanced rapidly in Europe, North America, and parts of Latin America, very little empirical research has been conducted in developing Southeast Asian economies, and specifically in Cambodia.

\hs Cambodia presents a uniquely compelling economic context for high-frequency digital price measurement:
\begin{enumerate}
    \item \textbf{High Dollarization and Currency Dualism:} Both the Cambodian Riel (KHR) and the US Dollar (USD) circulate widely in domestic commerce. Retailers regularly quote prices in either currency. Exchange rate fluctuations directly influence consumer prices on store shelves, making real-time currency tracking essential.
    \item \textbf{Extreme Expenditure Concentration:} As documented in the Cambodia Socio-Economic Survey (CSES 2023), food and non-alcoholic beverages account for \textbf{44.775\%} of the total national consumer basket. Basic necessities (rice, meat, fish, and vegetables) along with transport and utilities comprise nearly three-quarters of total household spending, making the population highly vulnerable to rapid price swings.
    \item \textbf{Rapid E-Commerce Expansion:} Over the past five years, the widespread adoption of smartphones, high-speed mobile internet, and the National Bank of Cambodia's digital payment system (Bakong) has catalyzed the growth of modern online supermarkets, digital pharmacies, and ride-hailing services.
\end{enumerate}

\hs Despite these unique conditions, national inflation tracking in Cambodia has continued to rely exclusively on monthly in-person market surveys. Prior to this research, no automated, end-to-end data engineering pipeline existed to continuously harvest, classify, calculate, and forecast daily consumer prices across all 12 UN COICOP divisions in Cambodia. This graduation thesis directly addresses this critical gap, providing both an academic contribution to applied econometric data science and a practical digital tool for national economic governance.

\section{Comparative Literature Matrix and Methodological Synthesis}
\hs Table~\ref{tab:literature_matrix} synthesizes the seminal academic publications, international statistical standards, and central bank studies that directly inform the four operational pillars of the Cambodia Daily CPI Medallion Pipeline, mapping each theoretical framework to its concrete implementation in the Cambodian system.

\begin{table}[H]
\centering
\scriptsize
\setlength{\tabcolsep}{3.5pt}
\begin{tabular}{|>{\raggedright\arraybackslash}p{2.5cm}|>{\centering\arraybackslash}p{1.7cm}|>{\raggedright\arraybackslash}p{3.3cm}|>{\raggedright\arraybackslash}p{3.2cm}|>{\raggedright\arraybackslash}p{3.5cm}|}
\hline
\rowcolor{blue!20}
\textbf{Paper / Study} & \textbf{Pipeline Pillar} & \textbf{Key Methodology \& Dataset} & \textbf{Reported Metric / Result} & \textbf{Implementation in Cambodia Pipeline} \\ \hline
\textbf{Billion Prices Project} \citep{cavallo2016billion} & Bronze Ingestion & Daily web crawlers, Jevons geometric mean, 22 countries & $>0.98$ correlation with official CPI; 3--14 day shock detection & Automated daily scraping of 20 retail sources across Cambodia \\ \hline
\textbf{Targeted Price Baskets} \citep{aparicio2021targeted} & Bronze Ingestion & Targeted online price baskets, high-frequency microdata sensing & Early detection of turning points and micro-price stickiness & Standardized daily basket tracking across all 12 UN COICOP divisions \\ \hline
\textbf{BERT COICOP NLP} \citep{berki2025nlp} & Silver AI Classification & Fine-tuned BERT transformer on scraped e-commerce titles & 94.56\% classification accuracy; 94.07\% precision & Multilingual vector embeddings mapping raw scraped listings into COICOP \\ \hline
\textbf{Project Spectrum} \citep{bis2024spectrum} & Silver AI Classification & Generative AI LLMs, dense embeddings, relational caching & $>95\%$ classification precision; latency $<3$ minutes & Tier 3 AI categorization with database vector caching for Cambodia \\ \hline
\textbf{CPI Manual} \citep{imf2020cpi} & Gold Axiomatic Index & Axiomatic index theory, elementary Jevons, missing price imputation & Global standard eliminating Carli upward drift (+1.18\%/yr) & Gold layer daily Jevons elementary index with CSES 2023 weights \\ \hline
\textbf{ML Inflation Forecast} \citep{medeiros2021forecasting} & Gold ML Forecasting & Random Forests, GBRT, LASSO with high-dimensional macro features & 10\%--25\% lower MSFE/RMSE over Phillips Curve and AR & Production LightGBM regressor predicting forward inflation \\ \hline
\textbf{Macro ML Forecasting} \citep{gouletcoulombe2022machine} & Gold ML Forecasting & GBRT, Random Forest, non-linear feature interaction analysis & Superiority of tree ensembles under economic shocks and holiday demand & Gradient boosted trees handling fuel resets and Cambodian holiday surges \\ \hline
\textbf{High-Freq Food Inflation} \citep{macias2023nowcasting} & Gold ML Forecasting & 2.4M daily food quotes, recursive time-series modeling & 15\%--35\% RMSE reduction over AR/ARIMA benchmarks & Food Division 01 7-day momentum computed as primary leading signal \\ \hline
\textbf{High-Freq ML Regress.} \citep{babii2022machine} & Gold ML Forecasting & Machine learning time-series regressions, walk-forward validation & Robust out-of-sample multi-horizon inflation forecasting & Enforced \texttt{TimeSeriesSplit} cross-validation across 7d, 14d, 30d horizons \\ \hline
\end{tabular}
\caption{Methodological Matrix of Literature Directly Informing the Cambodia Daily CPI Medallion Pipeline}
\label{tab:literature_matrix}
\end{table}

